\nonstopmode{}
\documentclass[letterpaper]{book}
\usepackage[times,inconsolata,hyper]{Rd}
\usepackage{makeidx}
\makeatletter\@ifl@t@r\fmtversion{2018/04/01}{}{\usepackage[utf8]{inputenc}}\makeatother
% \usepackage{graphicx} % @USE GRAPHICX@
\makeindex{}
\begin{document}
\chapter*{}
\begin{center}
{\textbf{\huge Package `robcdcc'}}
\par\bigskip{\large \today}
\end{center}
\ifthenelse{\boolean{Rd@use@hyper}}{\hypersetup{pdftitle = {robcdcc: Simulation and robust estimation for GARCH-cDCC process of order (1,1)}}}{}
\begin{description}
\raggedright{}
\item[Type]\AsIs{Package}
\item[Title]\AsIs{Simulation and robust estimation for GARCH-cDCC process of order
(1,1)}
\item[Version]\AsIs{1.0.2}
\item[Author]\AsIs{Eduardo}
\item[Maintainer]\AsIs{Eduardo }\email{eduardogabriel2077@gmail.com}\AsIs{}
\item[Description]\AsIs{Application with functions in Rcpp.}
\item[License]\AsIs{MIT + file LICENSE}
\item[Imports]\AsIs{Rcpp, truncnorm, dplyr, purrr, rlist, mvnfast}
\item[LinkingTo]\AsIs{Rcpp, RcppArmadillo}
\item[Suggests]\AsIs{testthat (>= 3.0.0)}
\item[Config/testthat/edition]\AsIs{3}
\item[RoxygenNote]\AsIs{7.2.0}
\item[NeedsCompilation]\AsIs{yes}
\item[SystemRequirements]\AsIs{C++17}
\end{description}
\Rdcontents{Contents}
\HeaderA{robcdcc-package}{Simulation and robust estimation for GARCH-cDCC process of order (1,1)}{robcdcc.Rdash.package}
\aliasA{robcdcc}{robcdcc-package}{robcdcc}
\keyword{package}{robcdcc-package}
%
\begin{Description}
Application with functions in Rcpp.
\end{Description}
%
\begin{Details}

The DESCRIPTION file:
This package was not yet installed at build time.\\{}

Index:  This package was not yet installed at build time.\\{}
\textasciitilde{}\textasciitilde{} An overview of how to use the package, including the most important functions \textasciitilde{}\textasciitilde{}
\end{Details}
%
\begin{Author}
Eduardo

Maintainer: Eduardo <eduardogabriel2077@gmail.com>
\end{Author}
%
\begin{References}
\textasciitilde{}\textasciitilde{} Literature or other references for background information \textasciitilde{}\textasciitilde{}
\end{References}
%
\begin{SeeAlso}
\textasciitilde{}\textasciitilde{} Optional links to other man pages, e.g. \textasciitilde{}\textasciitilde{}
\textasciitilde{}\textasciitilde{} \code{\LinkA{<pkg>}{<pkg>}} \textasciitilde{}\textasciitilde{}
\end{SeeAlso}
%
\begin{Examples}
\begin{ExampleCode}
~~ simple examples of the most important functions ~~
\end{ExampleCode}
\end{Examples}
\HeaderA{calc\_ht}{Calculate ht using GARCH}{calc.Rul.ht}
%
\begin{Description}
Calculate ht using GARCH
\end{Description}
%
\begin{Usage}
\begin{verbatim}
calc_ht(rt, eta)
\end{verbatim}
\end{Usage}
%
\begin{Arguments}
\begin{ldescription}
\item[\code{rt}] returns

\item[\code{eta}] GARCH parameters
\end{ldescription}
\end{Arguments}
%
\begin{Value}
volatility (ht)
\end{Value}
\HeaderA{calc\_robust\_devolatilized\_returns}{Calculate devolatilized returns using BIP-GARCH especification}{calc.Rul.robust.Rul.devolatilized.Rul.returns}
%
\begin{Description}
Calculate devolatilized returns using BIP-GARCH especification
\end{Description}
%
\begin{Usage}
\begin{verbatim}
calc_robust_devolatilized_returns(rt, eta_df, delta = 0.975)
\end{verbatim}
\end{Usage}
%
\begin{Arguments}
\begin{ldescription}
\item[\code{rt}] returns

\item[\code{eta\_df}] GARCH parameters

\item[\code{delta}] Quantile for Chi-Squared
\end{ldescription}
\end{Arguments}
%
\begin{Value}
devolatilized returns
\end{Value}
\HeaderA{estimateGARCH}{Calculate corretion factor}{estimateGARCH}
%
\begin{Description}
Calculate corretion factor
\end{Description}
%
\begin{Usage}
\begin{verbatim}
estimateGARCH(rt)
\end{verbatim}
\end{Usage}
%
\begin{Arguments}
\begin{ldescription}
\item[\code{rt}] return matrix
\end{ldescription}
\end{Arguments}
%
\begin{Value}
Fit results
\end{Value}
\HeaderA{fc}{Calculate correction factor}{fc}
%
\begin{Description}
Calculate correction factor
\end{Description}
%
\begin{Usage}
\begin{verbatim}
fc(delta, N)
\end{verbatim}
\end{Usage}
%
\begin{Arguments}
\begin{ldescription}
\item[\code{delta}] Quantile for Chi-Squared

\item[\code{N}] Dimension
\end{ldescription}
\end{Arguments}
%
\begin{Value}
Correction factor for the weighted variance estimator
\end{Value}
\HeaderA{RcppArmadillo-Functions}{Set of functions in example RcppArmadillo package}{RcppArmadillo.Rdash.Functions}
\aliasA{rcpparma\_bothproducts}{RcppArmadillo-Functions}{rcpparma.Rul.bothproducts}
\aliasA{rcpparma\_hello\_world}{RcppArmadillo-Functions}{rcpparma.Rul.hello.Rul.world}
\aliasA{rcpparma\_innerproduct}{RcppArmadillo-Functions}{rcpparma.Rul.innerproduct}
\aliasA{rcpparma\_outerproduct}{RcppArmadillo-Functions}{rcpparma.Rul.outerproduct}
%
\begin{Description}
These four functions are created when
\code{RcppArmadillo.package.skeleton()} is invoked to create a
skeleton packages.
\end{Description}
%
\begin{Usage}
\begin{verbatim}
rcpparma_hello_world()
rcpparma_outerproduct(x)
rcpparma_innerproduct(x)
rcpparma_bothproducts(x)
\end{verbatim}
\end{Usage}
%
\begin{Arguments}
\begin{ldescription}
\item[\code{x}] a numeric vector
\end{ldescription}
\end{Arguments}
%
\begin{Details}
These are example functions which should be largely
self-explanatory. Their main benefit is to demonstrate how to write a
function using the Armadillo C++ classes, and to have to such a
function accessible from R.
\end{Details}
%
\begin{Value}
\code{rcpparma\_hello\_world()} does not return a value, but displays a
message to the console.

\code{rcpparma\_outerproduct()} returns a numeric matrix computed as the
outer (vector) product of \code{x}.

\code{rcpparma\_innerproduct()} returns a double computer as the inner
(vector) product of \code{x}.

\code{rcpparma\_bothproducts()} returns a list with both the outer and
inner products.

\end{Value}
%
\begin{Author}
Dirk Eddelbuettel
\end{Author}
%
\begin{References}
See the documentation for Armadillo, and RcppArmadillo, for more details.
\end{References}
%
\begin{Examples}
\begin{ExampleCode}
  x <- sqrt(1:4)
  rcpparma_innerproduct(x)
  rcpparma_outerproduct(x)
\end{ExampleCode}
\end{Examples}
\HeaderA{robust\_calc\_ht}{Calculate ht using BIP-GARCH}{robust.Rul.calc.Rul.ht}
%
\begin{Description}
Calculate ht using BIP-GARCH
\end{Description}
%
\begin{Usage}
\begin{verbatim}
robust_calc_ht(rt, eta, delta)
\end{verbatim}
\end{Usage}
%
\begin{Arguments}
\begin{ldescription}
\item[\code{rt}] returns

\item[\code{eta}] GARCH parameters

\item[\code{delta}] Quantile for Chi-Squared
\end{ldescription}
\end{Arguments}
%
\begin{Value}
volatility (ht)
\end{Value}
\HeaderA{robust\_estimateGARCH}{Estimate GARCH(1,1) parameters using BIP-GARCH especification and  M-estimation}{robust.Rul.estimateGARCH}
%
\begin{Description}
Estimate GARCH(1,1) parameters using BIP-GARCH especification and 
M-estimation
\end{Description}
%
\begin{Usage}
\begin{verbatim}
robust_estimateGARCH(rt, cy, chisq, k = 30)
\end{verbatim}
\end{Usage}
%
\begin{Arguments}
\begin{ldescription}
\item[\code{rt}] returns

\item[\code{cy}] correction factor

\item[\code{chisq}] Quantile for Chi-Squared
\end{ldescription}
\end{Arguments}
%
\begin{Value}
devolatilized returns
\end{Value}
\HeaderA{simCDCC}{Simulate cDCC(1,1)}{simCDCC}
%
\begin{Description}
Simulate cDCC(1,1)
\end{Description}
%
\begin{Usage}
\begin{verbatim}
simCDCC(phi, eta, S, nobs, ndim, seed)
\end{verbatim}
\end{Usage}
%
\begin{Arguments}
\begin{ldescription}
\item[\code{phi}] Parameters for cDCC(1,1)

\item[\code{eta}] Parameters for GARCH(1,1)

\item[\code{S}] Intercept matrix

\item[\code{nobs}] Number of observations

\item[\code{ndim}] Number of dimensions

\item[\code{seed}] seed for random sampling
\end{ldescription}
\end{Arguments}
%
\begin{Value}
List with results
\end{Value}
\printindex{}
\end{document}
